%% kapa: exam=1 round=27 year=2016 period=2 form=A subject=accounting subject_ko=회계학 seq=7 number=47
%% source: Q-Net 감정평가사 기출문제 (cst00309) artlSeq=5169402 file=extracted/1008/5169402/제27회 감정평가사 1차 문제/제27회 감정평가사 2교시문제(A형).pdf
%% conversion: kordoc 4.24.51 -> tools/gen1.py (auto); review=proofread-sonnet; flags=none
\begin{question}{47}
무형자산에 관한 설명으로 옳지 않은 것은?
\begin{choices}
  \item 내부적으로 창출한 영업권은 자산으로 인식하지 않는다.
  \item 사업결합으로 인식하는 영업권은 사업결합에서 획득하였지만 개별적으로 식별하여 별도로 인식하는 것이 불가능한 그 밖의 자산에서 발생하는 미래경제적효익을 나타내는 자산이다.
  \item 무형자산을 창출하기 위한 내부 프로젝트를 연구단계와 개발단계로 구분할 수 없는 경우에는 그 프로젝트에서 발생한 지출은 모두 연구단계에서 발생한 것으로 본다.
  \item 자산에서 발생하는 미래경제적효익이 기업에 유입될 가능성이 높고 자산의 원가를 신뢰성 있게 측정할 수 있는 경우에만 무형자산을 인식한다.
  \item 경영자가 의도하는 방식으로 운용될 수 있으나 아직 사용하지 않고 있는 기간에 발생한 원가는 무형자산의 장부금액에 포함한다.
\end{choices}
\end{question}
\answer{5}
